% !TEX TS-program = xelatex
% !BIB TS-program = bibtex
\documentclass[12pt,letterpaper]{article}
\usepackage{style/dsc180reportstyle} % import dsc180reportstyle.sty
\newcommand{\lmul}{$\mathcal{L}$-Mul\xspace}

%%%%%%%%%%%%%%%%%%%%%%%%%%%%%%%%%%%%%%%%%%%%%%%%%%%%%%%%
%%%% Title and Authors
%%%%%%%%%%%%%%%%%%%%%%%%%%%%%%%%%%%%%%%%%%%%%%%%%%%%%%%%

\title{Insert Fancy Words About Hardware Acceleration Here}

\author{Kai Breese \\
  {\tt kbreese@ucsd.edu} \\\And
  Lukas Fullner \\
  {\tt lfullner@ucsd.edu} \\\And
  Justin Chou \\
  {\tt jtchou@ucsd.edu} \\\And
  Katelyn Abille \\
  {\tt kabille@ucsd.edu} \\\And
  Rajesh Gupta \\
  {\tt rgupta@ucsd.edu} \\}

\begin{document}
\maketitle

%%%%%%%%%%%%%%%%%%%%%%%%%%%%%%%%%%%%%%%%%%%%%%%%%%%%%%%%
%%%% Abstract and Links
%%%%%%%%%%%%%%%%%%%%%%%%%%%%%%%%%%%%%%%%%%%%%%%%%%%%%%%%

\begin{abstract}
    Efficient floating-point operations are a significant challenge in large neural networks and other computationally-intensive machine learning algorithms, where energy consumption and latency are key constraints. In this report, we present an implementation of the linear-complexity multiplication (\lmul) algorithm designed to approximate floating-point multiplication using addition \citep{luo2024addition}. By leveraging this approximation, \lmul achieves high precision with significantly lower computational cost compared to traditional methods of floating-point multiplication. Our approach aims to further deploy this method on FPGAs, offering a reduction in energy consumption and processing time. 
\end{abstract}

\begin{center}
% Website: \url{https://abc.github.io/} \\ % to  be added quarter 2
Code: \url{https://github.com/ninjakaib/hardware-accelerators}
\end{center}

\maketoc
\clearpage

%%%%%%%%%%%%%%%%%%%%%%%%%%%%%%%%%%%%%%%%%%%%%%%%%%%%%%%%
%%%% Main Contents
%%%%%%%%%%%%%%%%%%%%%%%%%%%%%%%%%%%%%%%%%%%%%%%%%%%%%%%%

\section{Introduction}

\section{Methods}

\section{Results}

\section{Discussion}

\section{Conclusion}

%%%%%%%%%%%%%%%%%%%%%%%%%%%%%%%%%%%%%%%%%%%%%%%%%%%%%%%%
%%%% Reference / Bibliography
%%%%%%%%%%%%%%%%%%%%%%%%%%%%%%%%%%%%%%%%%%%%%%%%%%%%%%%%

\makereference

\bibliography{reference}
\bibliographystyle{style/dsc180bibstyle}


%%%%%%%%%%%%%%%%%%%%%%%%%%%%%%%%%%%%%%%%%%%%%%%%%%%%%%%%
%%%% Appendix
%%%%%%%%%%%%%%%%%%%%%%%%%%%%%%%%%%%%%%%%%%%%%%%%%%%%%%%%

\clearpage
\makeappendix

\subsection{Training Details}

\subsection{Additional Figures}

\subsection{Additional Tables}


\end{document}